% Options for packages loaded elsewhere
\PassOptionsToPackage{unicode}{hyperref}
\PassOptionsToPackage{hyphens}{url}
\documentclass[
]{article}
\usepackage{xcolor}
\usepackage{amsmath,amssymb}
\setcounter{secnumdepth}{-\maxdimen} % remove section numbering
\usepackage{iftex}
\ifPDFTeX
  \usepackage[T1]{fontenc}
  \usepackage[utf8]{inputenc}
  \usepackage{textcomp} % provide euro and other symbols
\else % if luatex or xetex
  \usepackage{unicode-math} % this also loads fontspec
  \defaultfontfeatures{Scale=MatchLowercase}
  \defaultfontfeatures[\rmfamily]{Ligatures=TeX,Scale=1}
\fi
\usepackage{lmodern}
\ifPDFTeX\else
  % xetex/luatex font selection
\fi
% Use upquote if available, for straight quotes in verbatim environments
\IfFileExists{upquote.sty}{\usepackage{upquote}}{}
\IfFileExists{microtype.sty}{% use microtype if available
  \usepackage[]{microtype}
  \UseMicrotypeSet[protrusion]{basicmath} % disable protrusion for tt fonts
}{}
\makeatletter
\@ifundefined{KOMAClassName}{% if non-KOMA class
  \IfFileExists{parskip.sty}{%
    \usepackage{parskip}
  }{% else
    \setlength{\parindent}{0pt}
    \setlength{\parskip}{6pt plus 2pt minus 1pt}}
}{% if KOMA class
  \KOMAoptions{parskip=half}}
\makeatother
\usepackage{longtable,booktabs,array}
\usepackage{calc} % for calculating minipage widths
% Correct order of tables after \paragraph or \subparagraph
\usepackage{etoolbox}
\makeatletter
\patchcmd\longtable{\par}{\if@noskipsec\mbox{}\fi\par}{}{}
\makeatother
% Allow footnotes in longtable head/foot
\IfFileExists{footnotehyper.sty}{\usepackage{footnotehyper}}{\usepackage{footnote}}
\makesavenoteenv{longtable}
\setlength{\emergencystretch}{3em} % prevent overfull lines
\providecommand{\tightlist}{%
  \setlength{\itemsep}{0pt}\setlength{\parskip}{0pt}}
\usepackage{bookmark}
\IfFileExists{xurl.sty}{\usepackage{xurl}}{} % add URL line breaks if available
\urlstyle{same}
\hypersetup{
  hidelinks,
  pdfcreator={LaTeX via pandoc}}

\author{}
\date{}

\begin{document}

\section{\texorpdfstring{Hodge Structure and Gauge Equivalence of
\(\ell^1\) Defect
Fields}{Hodge Structure and Gauge Equivalence of \textbackslash ell\^{}1 Defect Fields}}\label{hodge-structure-and-gauge-equivalence-of-ell1-defect-fields}

\textbf{Author:} Jeremy H. Carroll \textbf{Date:} March 2026
\textbf{Version:} v4.0 (Pre-print)

\begin{center}\rule{0.5\linewidth}{0.5pt}\end{center}

\subsection{Classification of Results}\label{classification-of-results}

This paper separates rigorous mathematics from original constructions
and speculative analogies. Each result is annotated with its epistemic
layer:

\begin{longtable}[]{@{}
  >{\raggedright\arraybackslash}p{(\linewidth - 4\tabcolsep) * \real{0.2121}}
  >{\raggedright\arraybackslash}p{(\linewidth - 4\tabcolsep) * \real{0.2424}}
  >{\raggedright\arraybackslash}p{(\linewidth - 4\tabcolsep) * \real{0.5455}}@{}}
\toprule\noalign{}
\begin{minipage}[b]{\linewidth}\raggedright
Layer
\end{minipage} & \begin{minipage}[b]{\linewidth}\raggedright
Status
\end{minipage} & \begin{minipage}[b]{\linewidth}\raggedright
What it contains
\end{minipage} \\
\midrule\noalign{}
\endhead
\bottomrule\noalign{}
\endlastfoot
\textbf{A: Theorems} & Established mathematics & Standard results from
discrete algebraic topology and discrete \(\ell^2\) Hodge theory on
finite graphs without 2-cells {[}1, 7, 8{]}. \\
\textbf{B: Construction} & Original to this paper & The Cohomological
\(\ell^1\) Norm \(\Phi\) (terminology adopted here for the \(\ell^1\)
quotient norm on graph cohomology), the proof of its sparsity bias, LP
cut-duality, and the inflation quantification. \\
\textbf{C: Analogy} & Speculative & The structural parallel to
anisotropic total variation (graph-TV) denoising in image processing
{[}2, 13{]}. \\
\end{longtable}

\begin{center}\rule{0.5\linewidth}{0.5pt}\end{center}

\subsection{Abstract}\label{abstract}

The preceding works {[}14, 15, 16{]} establish that additive,
structure-preserving diagnostics on coboundary spaces uniquely induce an
\(\ell^1\) geometry under explicit axioms. Given this constraint, a
natural question is how to define a gauge-invariant measure of the
irreducible component of a defect field.

On finite graphs, classical discrete Hodge theory provides an
\(\ell^2\)-orthogonal decomposition separating exact and harmonic
components. However, this decomposition is not compatible with the
\(\ell^1\) metric: the \(\ell^2\)-orthogonal projection onto the
harmonic space does not, in general, minimize \(\ell^1\) cost within a
cohomology class.

We study the functional \[
\Phi([\delta]) = \min_{f \in C^0} \|\delta - d_0 f\|_1,
\] which is the standard \(\ell^1\) quotient norm on
\(C^1 / \operatorname{im}(d_0)\) (the minimum is attained; see
Proposition 8.2). While this optimization problem is classical
(appearing in total variation minimization and minimum homologous chain
problems {[}12{]}), we show that, under the structural constraints
identified in {[}14--16{]}, it is the unique gauge-invariant norm
compatible with the forced \(\ell^1\) geometry, in the sense that any
norm on \(H^1(G)\) induced from the \(\ell^1\) norm on \(C^1\) via the
quotient map must equal \(\Phi\) (Theorem 8.1, Remark 8.6).

We further relate \(\Phi\) to discrete Hodge decomposition, prove that
\(\ell^2\) projection generically inflates \(\ell^1\) cost (with an
explicit ratio of \(3 - 2/n\) on cycle graphs \(C_n\)), and characterize
\(\Phi\) via linear programming duality as the dual norm of bounded
divergence-free circulations. This provides a unified interpretation
connecting cohomology, \(\ell^1\) optimization, and network flows,
without introducing new optimization machinery.

\begin{center}\rule{0.5\linewidth}{0.5pt}\end{center}

\subsection{1. Introduction}\label{introduction}

The preceding papers in this series isolated the exact geometric
magnitude of interaction defects generated when correlated systems are
projected onto locally factorized manifolds. Paper 002 {[}16{]}
established that coordinate-wise additivity forces an \(\ell^1\)
geometry on finite graph cochains. Paper 001 {[}15{]} generalized this
to Banach presheaves via functoriality under bounded linear maps. Paper
000 {[}14{]} applied this geometric rigidity to show that projecting
dynamics onto a product-state manifold creates an \(\ell^1\)-measurable
curvature obstruction that persists at finite density.

A natural next question arises: given a defect field
\(\delta \in C^1(G)\), how should one decompose it into removable
(exact) and irremovable (cohomological) components? The classical tool
is the discrete Hodge decomposition, which operates in \(\ell^2\). But
since the preceding papers establish \(\ell^1\) as the uniquely forced
diagnostic geometry, a tension emerges.

Crucially, the preceding papers do not merely motivate the use of an
\(\ell^1\) metric --- they \emph{force} it. Throughout, we assume the
\(\ell^1\) geometry established in {[}14--16{]}; this paper studies its
consequences. Any decomposition or invariant associated with defect
fields must therefore be compatible with this geometry. The central
problem addressed here is thus not to choose a convenient norm, but to
determine the canonical cohomological structure induced by the uniquely
determined \(\ell^1\) metric.

\textbf{Classical Hodge decomposition does not generally hold in \(L^1\)
Banach spaces}, as they lack the inner-product structure necessary to
define orthogonal projections. We do not construct an \(\ell^1\) Hodge
decomposition. Instead, we use the classical \(\ell^2\) Hodge
decomposition as an algebraic coordinate system on \(C^1(G)\) and
evaluate its components under the \(\ell^1\) norm. By operating on
\textbf{finite graphs} (finite-dimensional discrete 1-complexes), the
vector space topology ensures that the \(\ell^2\) decomposition remains
algebraically available, while the resulting mismatch between orthogonal
projection in \(\ell^2\) and sparsity structure in \(\ell^1\) is made
precise.

We demonstrate that while the discrete \(\ell^2\) Hodge operators
dictate the irremovable topological \emph{shape} of the obstruction,
they generically increase its \(\ell^1\) \emph{magnitude}. To capture
the true physical penalty, we construct the \textbf{Cohomological
\(\ell^1\) Norm}: a well-defined functional
\(\Phi : H^1(G, \mathbb{R}) \to \mathbb{R}_{\geq 0}\) that evaluates the
topological defect not by \(\ell^2\) projection, but by infimal
\(\ell^1\) gauge cost.

\textbf{Scope and Relation to Existing Work.} We emphasize that the
optimization problem \(\min_f \|\delta - d_0 f\|_1\) is not new. It
appears in several established contexts, including anisotropic total
variation on graphs {[}2, 13{]}, \(\ell^1\) regression over incidence
matrices {[}4{]}, the minimum homologous chain problem {[}12{]}, and is
related to Gromov's bounded cohomology {[}11{]} in the
infinite-dimensional setting. Compressed sensing and basis pursuit
{[}5{]} exploit the same \(\ell^1\) sparsity structure, and shortest
cycle basis algorithms {[}6{]} address the underlying combinatorial
structure. The contribution of this paper is not to introduce a new
optimization framework, but to identify this functional as the unique
gauge-invariant norm compatible with the \(\ell^1\) geometry that is
structurally forced under the assumptions of {[}14--16{]}. Accordingly,
the results here should be understood as a structural interpretation and
organization of known optimization objects within a cohomological
framework, rather than as the introduction of new algorithms or
complexity results.

\textbf{Recall (Papers 001--002).} The \(\ell^1\) coboundary norm
\(I(s) = \sum_e \|\delta_e(s)\|\) is the unique (up to positive scalar)
edge-additive, faithful, functorial seminorm on the 1-cochain space
\(C^1(P, F)\) {[}15, 16{]}. In the presheaf language of {[}15{]}, the
coboundary is \(\delta_e(s) = r_{b,a}(s(b)) - s(a)\). When the Banach
presheaf \(F\) is the constant sheaf \(\mathbb{R}\) on a graph \(G\),
this reduces to \((d_0 f)(e) = f(v) - f(u)\), the total norm \(I(s)\)
becomes \(\|d_0 f\|_1\), and the quotient norm
\(\Phi([\delta]) = \min_f I(s - d_0 f)\) measures the irreducible cost
after gauge optimization. This paper develops the structure of \(\Phi\).

\textbf{Notation bridge.} The following dictionary connects the presheaf
notation of {[}14, 15{]} to the graph notation used in the remainder of
this paper:

\begin{longtable}[]{@{}
  >{\raggedright\arraybackslash}p{(\linewidth - 2\tabcolsep) * \real{0.5692}}
  >{\raggedright\arraybackslash}p{(\linewidth - 2\tabcolsep) * \real{0.4308}}@{}}
\toprule\noalign{}
\begin{minipage}[b]{\linewidth}\raggedright
Presheaf notation (Papers 000--001)
\end{minipage} & \begin{minipage}[b]{\linewidth}\raggedright
Graph notation (this paper)
\end{minipage} \\
\midrule\noalign{}
\endhead
\bottomrule\noalign{}
\endlastfoot
Poset \(P\), covering relations \(E_{\mathrm{cov}}(P)\) & Graph
\(G = (V, E)\) \\
Banach presheaf \(F\), stalks \(F(x)\), restrictions \(r_{b,a}\) &
Constant sheaf \(\mathbb{R}\), identity restrictions \\
Coboundary \(\delta_e(s) = r_{b,a}(s(b)) - s(a)\) & Discrete derivative
\((d_0 f)(e) = f(v) - f(u)\) \\
\(\ell^1\) coboundary norm \(I(s) = \sum_e \lVert\delta_e(s)\rVert\) &
\(\lVert d_0 f \rVert_1 = \sum_e \lvert (d_0 f)_e \rvert\) \\
\end{longtable}

\textbf{Conventions.} Throughout this paper, \(G = (V,E)\) denotes a
finite, connected, undirected graph. Each edge is assigned an arbitrary
but fixed orientation, used only for algebraic bookkeeping; all results
are independent of this choice. We work over the scalar field
\(\mathbb{R}\). Norms are denoted \(\lVert \cdot \rVert\) with
appropriate subscripts (\(\ell^1\), \(\ell^2\), \(\ell^\infty\));
absolute values of scalars are denoted \(|\cdot|\). The extension to
vector-valued defects \(\delta \in C^1(G, \mathbb{R}^k)\) proceeds
component-wise and is discussed in Section 10.

\begin{center}\rule{0.5\linewidth}{0.5pt}\end{center}

\subsection{\texorpdfstring{2. Gauge Equivalence \textbf{(Layer
A/B)}}{2. Gauge Equivalence (Layer A/B)}}\label{gauge-equivalence-layer-ab}

Before constructing the decomposition machinery, we introduce the
central physical notion that motivates the entire analysis.

A defect field \(\delta \in C^1(G)\) records the local projection error
across each edge. However, the absolute value of \(\delta\) on any
single edge is not physically meaningful: it depends on the choice of
local reference state at each vertex. A \textbf{gauge transformation} is
a reassignment of these local reference states, parameterized by a
vertex function \(f \in C^0(G) = \mathbb{R}^{|V|}\). We use ``gauge
equivalence'' in the sense of equivalence under reassignment of local
reference frames, analogous to the gauge freedom of electromagnetic
potentials \(A \mapsto A + d\Lambda\).

\textbf{Definition 2.1 (Gauge Equivalence).} Two defect fields
\(\delta, \delta' \in C^1(G)\) are \textbf{gauge equivalent} if their
difference is an exact gradient:
\[ \delta \sim \delta' \iff \delta - \delta' \in \operatorname{im}(d_0) \]

where \(d_0 : C^0 \to C^1\) is the discrete exterior derivative defined
in Section 3. The equivalence class
\([\delta] = \delta + \operatorname{im}(d_0)\) is the \textbf{gauge
class} of \(\delta\).

The physically meaningful quantity is not \(\delta\) itself but its
gauge class \([\delta]\). The natural gauge-invariant magnitude of an
obstruction is therefore:

\[ \operatorname{Cost}(\delta) = \min_{f \in C^0} \|\delta - d_0 f\|_1 \]

This is precisely the \(\ell^1\) quotient norm on
\(C^1 / \operatorname{im}(d_0)\). It represents the minimal \(\ell^1\)
cost achievable over all gauge-equivalent representatives. We will show
that this quantity defines a norm on the first cohomology group
\(H^1(G, \mathbb{R})\) and admits a clean dual formulation connecting it
to the topological cycle structure of \(G\).

\begin{center}\rule{0.5\linewidth}{0.5pt}\end{center}

\subsection{\texorpdfstring{3. The Discrete de Rham Complex
\textbf{(Layer
A)}}{3. The Discrete de Rham Complex (Layer A)}}\label{the-discrete-de-rham-complex-layer-a}

Let \(C^0 = \mathbb{R}^{|V|}\) denote the space of
\textbf{\(0\)-cochains} (vertex functions) and
\(C^1 = \mathbb{R}^{|E|}\) denote the space of \textbf{\(1\)-cochains}
(edge fields). Each edge \(e\) carries its fixed orientation
\(e = (u,v)\).

\textbf{Definition 3.1 (Discrete Exterior Derivative).} The operator
\(d_0 : C^0 \to C^1\) is the graph difference operator:
\[ (d_0 f)(e_{uv}) = f(v) - f(u) \]

Its matrix representation is the signed incidence matrix
\(B \in \mathbb{R}^{|E| \times |V|}\), with \(B_{e,v} = +1\),
\(B_{e,u} = -1\) for each oriented edge \(e = (u,v)\).

\textbf{Definition 3.2 (Discrete Codifferential / Divergence).} The
adjoint \(d_0^* : C^1 \to C^0\) is defined with respect to the standard
\(\ell^2\) inner products on \(C^0\) and \(C^1\):
\[ \langle d_0 f, \omega \rangle_{\ell^2} = \langle f, d_0^* \omega \rangle_{\ell^2} \]

Explicitly, \(d_0^* = B^T\) and the \textbf{divergence} of a
\(1\)-cochain \(\omega\) at a vertex \(u\) is the net signed flow:
\[ (d_0^* \omega)(u) = \sum_{e \text{ incoming at } u} \omega(e) - \sum_{e \text{ outgoing from } u} \omega(e) \]

The discrete de Rham complex for a graph without \(2\)-cells is:
\[ C^0 \xrightarrow{d_0} C^1 \xrightarrow{0} 0 \]

\begin{center}\rule{0.5\linewidth}{0.5pt}\end{center}

\subsection{\texorpdfstring{4. The Hodge Laplacian and Harmonic Cochains
\textbf{(Layer
A)}}{4. The Hodge Laplacian and Harmonic Cochains (Layer A)}}\label{the-hodge-laplacian-and-harmonic-cochains-layer-a}

Following Eckmann {[}1{]}, Horak--Jost {[}7{]}, and Lim {[}8{]}, the
full \(1\)-form Laplacian on a simplicial complex is:

\[ \Delta_1 = d_0 \, d_0^{*} + d_1^{*} \, d_1 \]

\textbf{Remark 4.1 (Convention and Missing 2-Cells).} For a graph (a
1-dimensional simplicial complex), there are no \(2\)-cells.
Consequently \(d_1 \equiv 0\), and the Laplacian collapses to:

\[ \Delta_1 = d_0 \, d_0^* \quad \text{(graphs without 2-cells)} \]

\textbf{Definition 4.2.} A \(1\)-cochain \(\omega \in C^1\) is
\textbf{harmonic} if \(\Delta_1 \omega = 0\):
\(\mathcal{H}^1(G) = \ker \Delta_1\).

\textbf{Proposition 4.3.} \emph{For a graph without \(2\)-cells,
\(\ker \Delta_1 = \ker d_0^*\).}

\emph{Proof.} \(\ker d_0^* \subseteq \ker \Delta_1\) is immediate.
Conversely, \(\Delta_1 \omega = 0\) implies
\(0 = \langle d_0 d_0^* \omega, \omega \rangle_{\ell^2} = \|d_0^* \omega\|_{\ell^2}^2\),
so \(d_0^* \omega = 0\). \(\blacksquare\)

\textbf{Proposition 4.4.} \emph{If \(G\) is connected,
\(\dim \mathcal{H}^1(G) = \beta_1(G) = |E| - |V| + 1\).}

\begin{center}\rule{0.5\linewidth}{0.5pt}\end{center}

\subsection{\texorpdfstring{5. The Discrete Hodge Decomposition
\textbf{(Layer
A)}}{5. The Discrete Hodge Decomposition (Layer A)}}\label{the-discrete-hodge-decomposition-layer-a}

On a general continuous manifold, the Hodge theorem splits any 1-form
into three orthogonal components:
\(\omega = d\alpha + \delta \beta + h\) (exact, coexact, and harmonic).
However, because our state space is a finite graph with no 2-cells,
\(d_1 = 0\), its adjoint \(d_1^* = 0\), and the coexact term vanishes
entirely. The decomposition reduces to exactly two functional pieces.

\textbf{Theorem 5.1 (Bipartite Hodge Decomposition).} \emph{Every
\(\delta \in C^1(G)\) decomposes algebraically uniquely as:}
\[ \delta = \delta_{\text{ex}} + \delta_{\text{harm}}, \quad \delta_{\text{ex}} \in \operatorname{im}(d_0), \quad \delta_{\text{harm}} \in \ker d_0^*, \quad \delta_{\text{ex}} \perp_{\ell^2} \delta_{\text{harm}} \]

\emph{Proof.} By the orthogonal complement theorem in finite dimensions,
\(C^1 = \operatorname{im}(d_0) \oplus (\operatorname{im}(d_0))^\perp\).
Since
\(\langle d_0 f, h \rangle_{\ell^2} = \langle f, d_0^* h \rangle_{\ell^2}\)
for all \(f, h\), we have
\((\operatorname{im}(d_0))^\perp = \ker(d_0^*)\). Therefore
\(C^1 = \operatorname{im}(d_0) \oplus \ker(d_0^*)\) with the two
summands \(\ell^2\)-orthogonal. \(\blacksquare\)

\textbf{Remark 5.2 (Finite-Dimensional Availability).} This theorem
explicitly relies on the finite-dimensionality of the graph.
Infinite-dimensional \(L^1\) spaces lack orthogonal complements, making
\(L^1\) Hodge theory notoriously pathological. We do not attempt an
\(\ell^1\) Hodge decomposition. Instead, we use the \(\ell^2\) Hodge
decomposition as an algebraic coordinate system on \(C^1(G)\) ---
guaranteed to exist in finite dimensions --- and then evaluate its
components under the \(\ell^1\) norm. The resulting mismatch reflects
the incompatibility between orthogonal projection in \(\ell^2\) and
sparsity structure in \(\ell^1\).

\textbf{Remark 5.3 (Topological Interpretation).} Under this framework,
\(\operatorname{im}(d_0)\) corresponds to potential difference
relationships (gauge transformations), while \(\ker d_0^*\) captures
network circulations. The harmonic component \(\delta_{\text{harm}}\) is
the gauge-invariant core, unchanged under
\(\delta \mapsto \delta - d_0 g\). For trees (\(\beta_1 = 0\)), all
defects are exact and trivially gauge-removable. For graphs with loops
(\(\beta_1 > 0\)), any non-zero harmonic component constitutes an
irreducible topological obstruction.

\begin{center}\rule{0.5\linewidth}{0.5pt}\end{center}

\subsection{\texorpdfstring{6. The \(\ell^1 / \ell^2\) Tension
\textbf{(Layer
A/B)}}{6. The \textbackslash ell\^{}1 / \textbackslash ell\^{}2 Tension (Layer A/B)}}\label{the-ell1-ell2-tension-layer-ab}

\textbf{Remark 6.1 (Layer A).} \(\|x + y\|_1 \leq \|x\|_1 + \|y\|_1\),
with equality iff \(x_e y_e \geq 0\) for all \(e\). This is the standard
triangle inequality.

\textbf{Remark 6.2 (Layer A).} Applying the triangle inequality to the
Hodge decomposition:
\(\|\delta\|_1 \leq \|\delta_{\text{ex}}\|_1 + \|\delta_{\text{harm}}\|_1\),
with equality iff there is edgewise sign compatibility between the exact
and harmonic components.

\textbf{Remark 6.3.} Generic sign conflicts make the inequality strict.
Thus, the \(\ell^2\) algebraic Hodge projection operator generically
yields strictly larger \(\ell^1\) component sums except in
sign-compatible cases.

\textbf{Proposition 6.4 (Failure of \(\ell^2\) Hodge Minimization for
\(\ell^1\) Cost) (Layer B).} Let \(\delta \in C^1(G)\) and let
\(\delta = \delta_{\mathrm{ex}} + \delta_{\mathrm{harm}}\) be its
\(\ell^2\) Hodge decomposition. Then:

\begin{enumerate}
\def\labelenumi{\arabic{enumi}.}
\item
  The harmonic component does not, in general, minimize \(\ell^1\) cost
  over the cohomology class \([\delta]\):
  \[\|\delta_{\mathrm{harm}}\|_1 \geq \Phi([\delta]),\] with strict
  inequality for generic \(\delta\).
\item
  Equality holds if and only if \(\delta_{\mathrm{harm}}\) is already an
  \(\ell^1\)-optimal representative of its cohomology class, which
  requires sign compatibility between the harmonic component and any
  exact correction.
\end{enumerate}

In particular, the \(\ell^2\)-orthogonal projection onto \(\ker d_0^*\)
is not an \(\ell^1\)-optimal representative of the cohomology class.

\emph{Proof.} By definition,
\(\Phi([\delta]) = \min_{f \in C^0} \|\delta - d_0 f\|_1 \leq \|\delta_{\mathrm{harm}}\|_1\)
since \(\delta_{\mathrm{harm}} = \delta - d_0 f_{\ell^2}\) for some
\(f_{\ell^2}\). Strict inequality holds whenever a different gauge
representative achieves lower \(\ell^1\) cost. By the crosspolytope
geometry of the \(\ell^1\) ball, the \(\ell^2\)-optimal representative
(which minimizes Euclidean distance to \(\ker d_0^*\)) generically
diffuses mass across coordinates, while the \(\ell^1\)-optimal
representative concentrates mass on minimal support. These coincide only
when the harmonic projection is already sparse (sign-compatible). The
examples below demonstrate strict inequality on all cycle graphs
\(C_n\), \(n \geq 3\). \(\blacksquare\)

\subsubsection{\texorpdfstring{6.1 Triangle Graph \textbf{(Layer
B)}}{6.1 Triangle Graph (Layer B)}}\label{triangle-graph-layer-b}

\textbf{Example 6.5.} On \(C_3\) with \(\beta_1 = 1\), cycle vector
\(c = (1,1,1)^T\). A sparse local defect \(\delta = (1,0,0)^T\) resolves
via \(\ell^2\) projection to: -
\(\delta_{\text{harm}} = \frac{1}{3}(1,1,1)^T\),
\(\quad \delta_{\text{ex}} = \frac{1}{3}(2,-1,-1)^T\) -
\(\|\delta_{\text{ex}}\|_1 + \|\delta_{\text{harm}}\|_1 = \frac{4}{3} + 1 = \frac{7}{3} > 1 = \|\delta\|_1\)
- Inflation ratio: \(7/3 \approx 2.33\)

\subsubsection{\texorpdfstring{6.2 Square Cycle \textbf{(Layer
B)}}{6.2 Square Cycle (Layer B)}}\label{square-cycle-layer-b}

\textbf{Example 6.6.} On \(C_4\) with \(\delta = (1,0,0,0)^T\): -
\(\delta_{\text{harm}} = \frac{1}{4}(1,1,1,1)^T\),
\(\quad \delta_{\text{ex}} = \frac{1}{4}(3,-1,-1,-1)^T\) -
\(\|\delta_{\text{harm}}\|_1 = 1\),
\(\quad \|\delta_{\text{ex}}\|_1 = \frac{3}{4} + \frac{1}{4} + \frac{1}{4} + \frac{1}{4} = \frac{6}{4} = \frac{3}{2}\)
-
\(\|\delta_{\text{ex}}\|_1 + \|\delta_{\text{harm}}\|_1 = \frac{5}{2} > 1 = \|\delta\|_1\)
- Inflation ratio: \(5/2 = 2.5\)

\subsubsection{\texorpdfstring{6.3 Inflation Ratio on Cycle Graphs
\textbf{(Layer
B)}}{6.3 Inflation Ratio on Cycle Graphs (Layer B)}}\label{inflation-ratio-on-cycle-graphs-layer-b}

\textbf{Proposition 6.7 (\(\ell^2\) Inflation Ratio on \(C_n\)).}
\emph{For a single-edge defect
\(\delta = \mathbf{e}_1 = (1, 0, \ldots, 0)^T\) on the cycle graph
\(C_n\) (\(n \geq 3\)), the \(\ell^2\) Hodge decomposition inflation
ratio is:}

\[ \frac{\|\delta_{\text{ex}}\|_1 + \|\delta_{\text{harm}}\|_1}{\|\delta\|_1} = 3 - \frac{2}{n} \]

\emph{Proof.} The harmonic space of \(C_n\) is
\(\ker(d_0^*) = \operatorname{span}(\mathbf{1})\), where
\(\mathbf{1} = (1,1,\ldots,1)^T\). The \(\ell^2\)-orthogonal projection
onto \(\ker(d_0^*)\) gives:
\[\delta_{\text{harm}} = \frac{\langle \delta, \mathbf{1} \rangle}{\langle \mathbf{1}, \mathbf{1} \rangle} \mathbf{1} = \frac{1}{n}\mathbf{1}\]
\[\delta_{\text{ex}} = \delta - \delta_{\text{harm}} = \left(\frac{n-1}{n}, \frac{-1}{n}, \ldots, \frac{-1}{n}\right)\]

Computing the \(\ell^1\) norms:
\[\|\delta_{\text{harm}}\|_1 = n \cdot \frac{1}{n} = 1\]
\[\|\delta_{\text{ex}}\|_1 = \frac{n-1}{n} + (n-1) \cdot \frac{1}{n} = \frac{2(n-1)}{n}\]

The inflation ratio is:
\[\frac{\|\delta_{\text{ex}}\|_1 + \|\delta_{\text{harm}}\|_1}{\|\delta\|_1} = \frac{\frac{2(n-1)}{n} + 1}{1} = \frac{2n - 2 + n}{n} = \frac{3n - 2}{n} = 3 - \frac{2}{n}\]

This grows monotonically with \(n\) and approaches 3 as
\(n \to \infty\). \(\blacksquare\)

\textbf{Remark 6.8 (Crosspolytope Bias).} There is a geometric reason
why the \(\ell^2\) harmonic projection generically increases \(\ell^1\)
magnitude. The mismatch stems from the fundamental difference between
Euclidean balls (\(\ell^2\)) and crosspolytopes (\(\ell^1\)) in high
dimensions. The \(\ell^2\) orthogonal projection minimizes Euclidean
distance but naturally diffuses mass across coordinates uniformly (e.g.,
\(1/3, 1/3, 1/3\)). The \(\ell^1\) cost, dictated by the crosspolytope
geometry, privileges sparse, extreme configurations (corners). When
mapping a sparse exact defect onto the harmonic cycle space, \(\ell^2\)
projection destroys its sparsity -- an inflation directly proportional
to the loss of corner-localization.

\begin{center}\rule{0.5\linewidth}{0.5pt}\end{center}

\subsection{\texorpdfstring{7. The Cohomological \(\ell^1\) Norm
\textbf{(Layer
B)}}{7. The Cohomological \textbackslash ell\^{}1 Norm (Layer B)}}\label{the-cohomological-ell1-norm-layer-b}

This section develops the \textbf{Cohomological \(\ell^1\) Norm}, a term
we adopt for the \(\ell^1\) quotient norm on
\(C^1/\operatorname{im}(d_0)\). The quotient norm construction is
standard in functional analysis {[}4, 9{]}; the \(\ell^1\) quotient norm
on homology classes appears in the bounded cohomology literature
{[}11{]} and the minimum homologous chain problem {[}12{]}. The
contribution here is not the optimization framework itself, but the
identification, under the additivity and functoriality assumptions
established in {[}14--16{]}, of \(\Phi\) as the unique gauge-invariant
norm on \(H^1(G)\) induced by the forced \(\ell^1\) geometry (Remark
8.6).

\subsubsection{\texorpdfstring{7.1 The \(\ell^1\) Quotient
Norm}{7.1 The \textbackslash ell\^{}1 Quotient Norm}}\label{the-ell1-quotient-norm}

As established in Section 2, the scalar magnitude invariant under
arbitrary local reference rescalings is the \textbf{\(\ell^1\) quotient
norm} on the space of cochains modulo exact gradients
\(C^1 / \operatorname{im}(d_0)\):
\[ \operatorname{Cost}(\delta) = \min_{f \in C^0} \|\delta - d_0 f\|_1 \]

\subsubsection{7.2 Descent to Cohomology}\label{descent-to-cohomology}

The gauge class \([\delta] = \delta + \operatorname{im}(d_0)\) contains
a unique element of \(\ker d_0^*\), namely
\(\delta_{\text{harm}} = \Pi_{\ker d_0^*}(\delta)\), the
\(\ell^2\)-orthogonal projection. This identification uses the
\(\ell^2\) Hodge decomposition to select a canonical representative
within each cohomology class; the \(\ell^1\) minimization itself is
independent of this choice. Therefore:

\[ \operatorname{Cost}(\delta) = \min_{r \in \delta_{\text{harm}} + \operatorname{im}(d_0)} \|r\|_1 \]

This evaluates to a purely cohomological constraint. It depends only on
the cohomology class \([\delta]\), with the harmonic representative
providing a canonical coordinate in \(H^1(G, \mathbb{R})\).

\subsubsection{7.3 The Cohomological Norm}\label{the-cohomological-norm}

\textbf{Proposition 7.1 (Existence of Minimizers).} \emph{The minimum
\(\min_{f \in C^0} \|\delta - d_0 f\|_1\) is attained for every
\(\delta \in C^1(G)\).}

\emph{Proof.} The objective \(f \mapsto \|\delta - d_0 f\|_1\) is
continuous and coercive on \(C^0 / \ker(d_0) \cong \mathbb{R}^{|V|-1}\)
(since \(\ker(d_0) = \operatorname{span}(\mathbf{1})\) for connected
\(G\) and \(\|d_0 f\|_1 \to \infty\) as \(\|f\|_2 \to \infty\) on this
quotient). A continuous coercive function on a finite-dimensional space
attains its infimum {[}9{]}. \(\blacksquare\)

\emph{Every minimizer \(r^* = \delta - d_0 f^*\) satisfies the
first-order optimality condition \(B^T s^* = 0\) where
\(s^* \in \partial\|\cdot\|_1(r^*)\) is a subgradient {[}9, Section
23{]}. That is, \(s^*\) is a divergence-free flow bounded entrywise by
1.}

\textbf{Theorem 7.2 (Cohomological \(\ell^1\) Norm).} \emph{The
functional}

\[ \Phi : H^1(G, \mathbb{R}) \to \mathbb{R}_{\geq 0}, \qquad \Phi([\delta]) = \min_{r \in [\delta]} \|r\|_1 \]

\emph{is a norm on \(H^1(G, \mathbb{R})\).}

\emph{Proof.} \textbf{Positive definiteness.} If \(\Phi([\delta]) = 0\),
then by Proposition 7.1 the minimum is attained: there exists \(f^*\)
with \(\|\delta - d_0 f^*\|_1 = 0\), hence \(\delta = d_0 f^*\), so
\([\delta] = 0\). Conversely, \([\delta] = 0\) gives
\(\delta \in \operatorname{im}(d_0)\), so \(\Phi([\delta]) = 0\).

\textbf{Absolute homogeneity.} For \(\lambda \neq 0\):
\(\Phi([\lambda \delta]) = \min_f \|\lambda \delta - d_0 f\|_1 = \min_f \|\lambda(\delta - d_0(f/\lambda))\|_1 = |\lambda| \min_g \|\delta - d_0 g\|_1 = |\lambda| \Phi([\delta])\),
where \(g = f/\lambda\). For \(\lambda = 0\): \(\Phi([0]) = 0\).

\textbf{Triangle inequality.} Let \(f^*, g^*\) be minimizers for
\([\delta]\) and \([\delta']\) respectively (existing by Proposition
7.1). Choosing \(f = f^* + g^*\) in the minimization for
\([\delta + \delta']\):
\[\Phi([\delta + \delta']) = \min_f \|\delta + \delta' - d_0 f\|_1 \leq \|(\delta - d_0 f^*) + (\delta' - d_0 g^*)\|_1 \leq \|\delta - d_0 f^*\|_1 + \|\delta' - d_0 g^*\|_1 = \Phi([\delta]) + \Phi([\delta'])\]
\(\blacksquare\)

\subsubsection{7.4 Sparsity Bias}\label{sparsity-bias}

\textbf{Remark 7.3 (Geometric Sparsity Bias).} \(\ell^1\) minimization
inherently contacts the bounding faces of the crosspolytope geometry on
low-dimensional faces (minimizing support). Consequently, optimization
concentrates defects onto minimal edge cycles rather than expanding them
as uniform fractional obstructions like the \(\ell^2\) harmonic
projection inevitably does.

\subsubsection{7.5 Structure on Cycle Bases and Minimum Weight
Combinations}\label{structure-on-cycle-bases-and-minimum-weight-combinations}

While minimizing \(\ell^1\) cost generally necessitates linear
programming, the structural geometry of the cycle space allows for a
direct evaluation on the graph's loops under favorable conditions.

\textbf{Theorem 7.4 (Evaluation on Edge-Disjoint Cycles).} \emph{Let
\(\{C_1, \ldots, C_{\beta_1}\}\) be a cycle basis for \(H^1(G)\) in
which the cycles are pairwise edge-disjoint and vertex-disjoint (i.e.,
no two cycles share an edge or a vertex). Given the harmonic expansion
\(\delta_{\text{harm}} = \sum_{i=1}^{\beta_1} a_i C_i\), the
cohomological norm is:}

\[ \Phi([\delta]) = \sum_{i=1}^{\beta_1} |a_i| \|C_i\|_1 \]

\emph{Proof.} Since the cycles are vertex-disjoint, their vertex sets
\(V_i\) are pairwise disjoint. The potential \(f \in C^0\) decomposes as
\(f = \sum_i f_i\) where \(f_i\) is supported on \(V_i\). The exact form
\(d_0 f\) restricted to edge set \(E_i\) depends only on \(f_i\)
(vertices outside \(V_i\) contribute zero to edges in \(E_i\)).
Therefore:

\[\min_f \|\delta - d_0 f\|_1 = \sum_{i=1}^{\beta_1} \min_{f_i \in \mathbb{R}^{V_i}} \sum_{e \in E_i} |a_i (C_i)_e - (d_0 f_i)_e|\]

On each component, the cohomology class \([a_i C_i]\) is
one-dimensional, and any exact correction on a single cycle changes the
representative by \(d_0 f_i\) for \(f_i\) supported on \(V_i\). Since
\(C_i\) is a single cycle, the minimum \(\ell^1\) representative of the
class \([a_i C_i]\) has cost \(|a_i| \|C_i\|_1\) (any exact correction
on a cycle merely redistributes, never reduces, the total absolute
circulation). \(\blacksquare\)

\textbf{Remark 7.5.} When cycles share vertices but not edges, the
decomposition may still hold in practice, but the proof requires a more
delicate LP analysis showing that the shared-vertex coupling constraints
are independently satisfiable. When cycles share edges, the general case
requires the linear program to resolve interactions, corresponding to
the \textbf{minimum weight homologous chain} {[}12{]}.

\textbf{Remark 7.6 (Geometry of the Cycle Polytope).} Under the \(\Phi\)
functional, the metric geometry of \(H^1(G)\) behaves structurally as a
polyhedral metric space whose unit locus mirrors the network's bounded
cycle polytope
(\(\operatorname{conv}\{\pm C_1, \ldots, \pm C_{\beta_1}\}\)).

\textbf{Remark 7.7 (Uniqueness of \(\Phi\)).} The norm \(\Phi\) is the
unique norm on \(H^1(G)\) induced via the \(\ell^1\) quotient
construction on \(C^1(G)\) through the quotient map
\(C^1 \to C^1/\operatorname{im}(d_0)\). This follows from the standard
construction: the quotient norm \(\|[x]\| = \inf\{\|y\| : y \in [x]\}\)
is the largest norm on the quotient space making the projection
non-expansive {[}9, Section 2{]}. Since {[}14--16{]} establish
\(\ell^1\) as the uniquely forced norm on \(C^1\), the quotient
construction uniquely determines \(\Phi\). Choosing any other
gauge-invariant norm on \(H^1\) (e.g., the \(\ell^2\) norm on harmonic
representatives) would break compatibility with the induced \(\ell^1\)
geometry.

\begin{center}\rule{0.5\linewidth}{0.5pt}\end{center}

\subsection{\texorpdfstring{8. Total Variation Denoising and Fast
Algorithms \textbf{(Layer
C)}}{8. Total Variation Denoising and Fast Algorithms (Layer C)}}\label{total-variation-denoising-and-fast-algorithms-layer-c}

The gauge optimization \(\min_f \|\delta - B f\|_1\) exhibits structural
isomorphism to \textbf{anisotropic total variation (TV) denoising}
extensively developed in image filtering {[}2, 13{]}. Although the
optimization is formally identical to graph total variation
minimization, its role here is different: it arises as the quotient norm
induced by the uniquely determined \(\ell^1\) geometry on coboundary
spaces {[}14--16{]}, rather than as a regularization principle or signal
recovery objective. Setting the mapping between computational models:

\begin{longtable}[]{@{}ll@{}}
\toprule\noalign{}
Imaging & Defect Topology \\
\midrule\noalign{}
\endhead
\bottomrule\noalign{}
\endlastfoot
Image Signal & Vertex Potential (\(f\)) \\
Image Gradient (\(\nabla\)) & Exterior Derivative (\(d_0 = B\)) \\
Observed Gradient (\(y\)) & Source Defect Field (\(\delta\)) \\
TV Residual (\(y - \nabla u\)) & Topological Obstruction (\(r\)) \\
\end{longtable}

Consequently, the optimization objective directly mirrors \textbf{graph
total variation (graph-TV)}. Where classical TV identifies structural
borders (sparse semantic image edges), the framework localizes purely
topological features (\textbf{sparse cycle contradictions}).

\textbf{Computational Remark.} This overlap permits the use of
algorithms designed for large-scale imaging tasks. Rather than utilizing
generic \(O(|E|^3)\) interior point bounds, the \(\ell^1\) quotient norm
is computable using splitting protocols (Chambolle--Pock primal-dual
{[}2{]}, Split Bregman, ADMM), which admit per-iteration complexity
scaling linearly in \(|E|\). Whether these methods achieve comparable
convergence rates in the graph cohomology setting is an empirical
question not addressed here; we note only the formal structural
correspondence.

\begin{center}\rule{0.5\linewidth}{0.5pt}\end{center}

\subsection{\texorpdfstring{9. Cut Duality and Flow Circulation
\textbf{(Layer
B)}}{9. Cut Duality and Flow Circulation (Layer B)}}\label{cut-duality-and-flow-circulation-layer-b}

We establish that the topological magnitude of the norm is strictly
characterized by valid internal network configurations.

\textbf{Theorem 9.1 (Cut Duality).}
\[ \Phi([\delta]) = \max_{\phi} \left\{ \langle \phi, \delta \rangle : \|\phi\|_\infty \leq 1, \; d_0^* \phi = 0 \right\} \]

\emph{Proof.} The primal problem \(\min_{f} \|\delta - B f\|_1\) is
reformulated as a linear program by introducing auxiliary variables
\(t_e \geq 0\):

\[\text{Primal: } \min_{f, t} \sum_e t_e \quad \text{subject to} \quad -t_e \leq \delta_e - (Bf)_e \leq t_e, \quad t_e \geq 0\]

Let \(\phi^+_e, \phi^-_e \geq 0\) be the dual variables for the upper
and lower bound constraints respectively. Setting
\(\phi_e = \phi^+_e - \phi^-_e\), the dual problem becomes:

\[\text{Dual: } \max_{\phi} \langle \phi, \delta \rangle \quad \text{subject to} \quad B^T \phi = 0, \quad |\phi_e| \leq 1 \text{ for all } e\]

The constraint \(B^T \phi = d_0^* \phi = 0\) requires \(\phi\) to be
divergence-free (a circulation). The bound \(|\phi_e| \leq 1\)
constrains \(\phi \in B_\infty\). Strong duality holds because the
problem is a finite-dimensional LP with a non-empty primal feasible set
(\(f = 0\) is feasible) and a bounded primal objective {[}4, Chapter
5{]}. \(\blacksquare\)

\textbf{Remark 9.2 (Dual Norm of Bounded Circulations).} The
\(d_0^* \phi = 0\) condition fixes \(\phi \in Z_1(G)\), the cycle space.
The bounded infinity constraint yields:

\[ \Phi([\delta]) = \max_{\phi \in Z_1(G) \cap B_\infty} \langle \phi, \delta \rangle \]

Hence, \(\Phi\) resolves as the exact \textbf{dual norm of bounded
circulations}. Within the network interpretation, it isolates the
maximal potential flow pairing compatible with the measured
contradiction field \(\delta\).

\begin{center}\rule{0.5\linewidth}{0.5pt}\end{center}

\subsection{\texorpdfstring{10. Extension to Vector-Valued Defects
\textbf{(Layer
A/B)}}{10. Extension to Vector-Valued Defects (Layer A/B)}}\label{extension-to-vector-valued-defects-layer-ab}

For \(\delta \in C^1(G, \mathbb{R}^k)\), define the separable \(\ell^1\)
norm \(\|\delta\|_1 = \sum_{e \in E} \sum_{j=1}^k |\delta_{e,j}|\)
(i.e., \(\ell^1\) over edges, then \(\ell^1\) over components). Under
this norm, the decomposition and cohomological norm extend
component-wise: \(\Phi([\delta]) = \sum_{j=1}^k \Phi([\delta^{(j)}])\).

\begin{center}\rule{0.5\linewidth}{0.5pt}\end{center}

\subsection{11. Discussion}\label{discussion}

\subsubsection{11.1 What This Paper
Establishes}\label{what-this-paper-establishes}

The discrete \(\ell^2\) Hodge decomposition canonically identifies the
topological shape of a cohomology class, but it does not respect the
\(\ell^1\) geometry established in the preceding papers {[}14, 15, 16{]}
as the uniquely forced diagnostic norm. The Cohomological \(\ell^1\)
Norm \(\Phi([\delta])\) resolves this by optimizing over
gauge-equivalent representatives to find the one of minimal \(\ell^1\)
cost.

The conceptual progression is: 1. \textbf{Topological classification}
(Theorem 5.1, Layer A): Finite dimensions rescue the algebraic Hodge
decomposition, categorizing gauge-removable and irreducible
obstructions. 2. \textbf{\(\ell^1/\ell^2\) tension} (Section 6, Layer
B): Evaluates Euclidean vs crosspolytope geometry, quantifying the
harmonic \(\ell^2\) projection cost inflation at a ratio of \(3 - 2/n\)
on cycle graphs. 3. \textbf{Cohomological norm} (Theorem 7.2, Layer B):
Formalizes \(\Phi([\delta]) = \min_r \|r\|_1\) as the \(\ell^1\)
quotient norm evaluating the defect's cycle space structure. 4.
\textbf{Circulation duality} (Theorem 9.1, Layer B): Expresses
obstruction severity as the maximal bounded divergence-free flow
pairing. 5. \textbf{Image processing parallelism} (Section 8, Layer C):
Notes the structural correspondence to TV denoising methods.

\subsubsection{11.2 Relation to Prior
Work}\label{relation-to-prior-work}

The discrete de Rham complex on graphs is due to Eckmann {[}1{]}. The
bipartite Hodge decomposition on graphs without 2-cells is standard (see
Lim {[}8{]} for a modern treatment). The LP duality formulation of
\(\Phi\) connects to classical network flow theory (Ford--Fulkerson
{[}3{]}) and the coarea formula for graph total variation (Chambolle et
al.~{[}2{]}). The \(\ell^1\) quotient norm on cohomology is related to
Gromov's bounded cohomology {[}11{]}, though the finite-dimensional
setting here avoids the subtleties of the infinite-dimensional theory.
The quotient norm construction itself is standard in functional analysis
{[}9{]}. The discrete calculus framework is developed systematically in
{[}10{]}.

\subsubsection{11.3 Position in the Unified
Program}\label{position-in-the-unified-program}

This paper completes the structural program initiated in the preceding
works by resolving the final open question: given that \(\ell^1\) is the
uniquely forced geometry on coboundary defects, what is the canonical
gauge-invariant decomposition and measurement of such defects?

The progression is now complete:

\begin{enumerate}
\def\labelenumi{\arabic{enumi}.}
\tightlist
\item
  \textbf{Combinatorial Level (Paper 002)} {[}16{]}: Local additivity
  over disjoint coordinate supports and edge separability force
  \(\ell^1\) geometry on finite cochains.
\item
  \textbf{Functional Level (Paper 001)} {[}15{]}: Structural
  compatibility (functoriality under bounded linear maps) extends this
  rigidity to Banach presheaves.
\item
  \textbf{Dynamical Level (Paper 000)} {[}14{]}: Projection of linear
  dynamics onto nonlinear manifolds generates intrinsic coboundary
  defects measured in this geometry, persisting at non-vanishing
  density.
\item
  \textbf{Cohomological Level (this paper)}: The forced \(\ell^1\)
  geometry induces a canonical quotient norm \(\Phi\) on \(H^1\),
  providing the unique gauge-invariant measure of irreducible
  obstruction (Remark 7.7). The \(\ell^2\) Hodge decomposition
  identifies the topological shape; \(\Phi\) measures the true
  \(\ell^1\) cost.
\end{enumerate}

Thus, the \(\ell^1\) structure is not an isolated phenomenon but
persists across combinatorial, functional, dynamical, and topological
layers. The measurement, decomposition, and optimization of defect
fields are not independent choices, but are constrained by a single
underlying geometric requirement.

\emph{Remark (Observer-theoretic grounding).} The axioms forcing
\(\ell^1\) geometry at each level of this escalation ladder admit an
interpretation in terms of observational consistency: they correspond to
the minimal structural conditions under which a local agent can
consistently detect and aggregate defects. Edge locality requires local
evaluability; coordinate additivity requires that independent channels
contribute without cancellation; symmetry requires indifference to
relabeling of structurally identical subsystems. Under this reading, the
four-level program --- from combinatorial cochains through Banach
presheaves and dynamics to the cohomological quotient norm developed
here --- admits an interpretation as grounded in a single operational
criterion: the existence of consistent local observers. The
gauge-invariant norm \(\Phi\) inherits this interpretation: it measures
the irreducible obstruction as seen by any observer who satisfies these
consistency requirements. This perspective is developed in {[}14,
Section 1.2{]} and parallels the combinatorial and functional-analytic
versions established in {[}16, Section 3.3{]} and {[}15, Section 6.4{]}.

\subsubsection{11.4 Limitations and Future
Work}\label{limitations-and-future-work}

The inflation analysis in Section 6 is restricted to cycle graphs
\(C_n\) with a single-edge localized defect. A complete characterization
of the inflation factor for general graphs remains open. The TV
denoising analogy in Section 8 is structural rather than rigorous;
formalizing the correspondence and establishing convergence rates is
left to future work. The restriction to graphs without 2-cells means the
coexact term vanishes. Extending the analysis to simplicial 2-complexes,
where the full tripartite Hodge decomposition
\(\delta = d\alpha + \delta\beta + h\) is active, is a natural next
step.

\subsubsection{11.5 Closing Remark}\label{closing-remark}

The results of this paper do not introduce new optimization problems or
algorithms; rather, they identify the structural role of an existing
\(\ell^1\) minimization within a cohomological framework. When combined
with the preceding classification results {[}14--16{]}, this shows that
the measurement, decomposition, and optimization of defect fields are
not independent choices, but are constrained by a single underlying
geometric requirement. The \(\ell^1\) quotient norm \(\Phi\) thus serves
as the consistent endpoint of this structure, linking local defect
aggregation, global topology, and convex optimization in a common
framework.

\begin{center}\rule{0.5\linewidth}{0.5pt}\end{center}

\subsection{Acknowledgments}\label{acknowledgments}

No external funding. No conflicts of interest.

\begin{center}\rule{0.5\linewidth}{0.5pt}\end{center}

\subsection{References}\label{references}

{[}1{]} B. Eckmann, ``Harmonische Funktionen und Randwertaufgaben in
einem Komplex,'' \emph{Commentarii Mathematici Helvetici} \textbf{17}
(1944/45), 240--255.

{[}2{]} A. Chambolle and T. Pock, ``A first-order primal-dual algorithm
for convex problems with applications to imaging,'' \emph{J. Math.
Imaging Vision} \textbf{40} (2011), 120--145.

{[}3{]} L. R. Ford and D. R. Fulkerson, \emph{Flows in Networks},
Princeton University Press, 1962.

{[}4{]} S. Boyd and L. Vandenberghe, \emph{Convex Optimization},
Cambridge University Press, 2004.

{[}5{]} E. J. Cand`es and T. Tao, ``Decoding by linear programming,''
\emph{IEEE Trans. Inform. Theory} \textbf{51} (2005), 4203--4215.

{[}6{]} J. D. Horton, ``A polynomial-time algorithm to find the shortest
cycle basis of a graph,'' \emph{SIAM J. Comput.} \textbf{16} (1987),
358--366.

{[}7{]} D. Horak and J. Jost, ``Spectra of combinatorial Laplace
operators on simplicial complexes,'' \emph{Adv. Math.} \textbf{244}
(2013), 303--336.

{[}8{]} L.-H. Lim, ``Hodge Laplacians on Graphs,'' \emph{SIAM Review}
\textbf{62} (2020), 685--715.

{[}9{]} R. T. Rockafellar, \emph{Convex Analysis}, Princeton University
Press, 1970.

{[}10{]} R. Grady and A. Polimeni, \emph{Discrete Calculus}, Springer,
2010.

{[}11{]} M. Gromov, ``Volume and bounded cohomology,''
\emph{Publications Math\textquotesingle ematiques de
l'IH\textquotesingle ES} \textbf{56} (1982), 5--99.

{[}12{]} T. K. Dey, A. N. Hirani, and B. Krishnamoorthy, ``Optimal
homologous cycles, total unimodularity, and linear programming,''
\emph{SIAM J. Comput.} \textbf{40} (2011), 1026--1044.

{[}13{]} L. I. Rudin, S. Osher, and E. Fatemi, ``Nonlinear total
variation based noise removal algorithms,'' \emph{Physica D: Nonlinear
Phenomena} \textbf{60} (1992), 259--268.

{[}14{]} J. H. Carroll, ``Single Obstruction Theory: Retraction
Nonlinearity, \(\ell^1\) Rigidity, and Density Scaling,'' Zenodo, 2026.
https://doi.org/10.5281/zenodo.19151803

{[}15{]} J. H. Carroll, ``The Free \(\ell^1\) Seminorm on Banach
Presheaf Coboundaries,'' Zenodo, 2026.
https://doi.org/10.5281/zenodo.19152115

{[}16{]} J. H. Carroll, ``Coordinate-Wise Additivity and the \(\ell^1\)
Norm on Finite Graph Cochains,'' Zenodo, 2026.
https://doi.org/10.5281/zenodo.18945342

{[}17{]} S. Kakutani, ``Concrete representation of abstract
\((L)\)-spaces and the mean ergodic theorem,'' \emph{Annals of
Mathematics} \textbf{42} (1941), 523--537.

\end{document}
